% 简单函数积分的“高度×测度”：[0,3]上 s=2 在 A1=[0,1)∪[2,2.5)，s=0.5 在 A2=[1,2)∪[2.5,3]；λ(A1)=λ(A2)=1.5，积分=2·1.5+0.5·1.5=3.75。
\begin{tikzpicture}[x=1.7cm,y=0.85cm,font=\fontsize{8}{10}\selectfont,
  ax/.style={draw=FigureStroke,line width=.75pt,-{Stealth[length=4pt]}},
  hi/.style={draw=FigureBlue,fill=FigureBlueFill,line width=.8pt},
  lo/.style={draw=FigureGreen,fill=FigureGreenFill,line width=.8pt},
  note/.style={font=\FigureNoteFont\fontsize{7.5}{9}\selectfont,text=FigureMuted}]
  \draw[hi] (0,0) rectangle (1,2);
  \draw[hi] (2,0) rectangle (2.5,2);
  \draw[lo] (1,0) rectangle (2,0.5);
  \draw[lo] (2.5,0) rectangle (3,0.5);
  \draw[ax] (-0.1,0)--(3.35,0) node[right,text=FigureInk]{$x$};
  \draw[ax] (0,-0.1)--(0,2.6) node[above,text=FigureInk]{$s(x)$};
  \foreach \y/\t in {0.5/0.5,2/2}{\draw[FigureStroke,line width=.6pt] (-0.05,\y)--(0,\y) node[left=2pt,text=FigureInk]{$\t$};}
  \foreach \x in {0,1,2,2.5,3}{\draw[FigureStroke,line width=.6pt] (\x,0)--(\x,-0.08) node[below,text=FigureInk]{$\x$};}
  \node[text=FigureBlue] at (0.5,1.2) {$c_1=2$};
  \node[text=FigureBlue] at (2.25,2.3) {$c_1$};
  \node[text=FigureGreen] at (1.5,0.85) {$c_2=0.5$};
  \node[text=FigureGreen] at (2.75,0.85) {$c_2$};
  % 集合标记
  \draw[FigureBlue,line width=1.6pt] (0,-0.72)--(1,-0.72) (2,-0.72)--(2.5,-0.72);
  \draw[FigureGreen,line width=1.6pt] (1,-1.25)--(2,-1.25) (2.5,-1.25)--(3,-1.25);
  \node[anchor=west,text=FigureBlue] at (3.08,-0.72) {$A_1$，$\lambda(A_1)=1.5$};
  \node[anchor=west,text=FigureGreen] at (3.08,-1.25) {$A_2$，$\lambda(A_2)=1.5$};
  \node[note,anchor=west] at (3.08,1.4) {同一高度的块合并成一层};
  \node[note,anchor=west] at (3.08,1.0) {积分$=2\times1.5+0.5\times1.5=3.75$};
\end{tikzpicture}
